\documentclass[a4paper]{article}

%------------------------------------------------
% الحزم
%------------------------------------------------

\usepackage{enumitem}
\usepackage[margin=2.5cm]{geometry}
\usepackage{amsmath,amsfonts,amssymb}
\usepackage[charter]{mathdesign}
\usepackage{graphicx}
\usepackage{array}
\usepackage{xcolor}
\usepackage{needspace}
\usepackage{fancyhdr}
\usepackage{polyglossia}

%------------------------------------------------
% اللغات والخطوط
%------------------------------------------------

\setdefaultlanguage[
    calendar=gregorian,
    numerals=maghrib
]{arabic}

\setotherlanguage{english}

% إذا لم تكن الخطوط مثبتة في الجهاز، تُؤخذ من المجلد ../../fonts/
\IfFontExistsTF{Amiri}{%
    \newfontfamily\arabicfont[Script=Arabic, Scale=1.2]{Amiri}
}{%
    \newfontfamily\arabicfont[
        Script=Arabic, Scale=1.2,
        Path=../../fonts/, Extension=.ttf,
        UprightFont=*-Regular, BoldFont=*-Bold,
        ItalicFont=*-Italic, BoldItalicFont=*-BoldItalic
    ]{Amiri}
}

\IfFontExistsTF{AlBattar}{%
    \newfontfamily\arabicfontsf[Script=Arabic, Scale=1.3]{AlBattar}
    \newfontfamily\khatt[Script=Arabic, Scale=1.3]{AlBattar}
}{%
    \newfontfamily\arabicfontsf[
        Script=Arabic, Scale=1.3, Path=../../fonts/,
        BoldFont=ae_AlBattar.ttf
    ]{ae_AlBattar.ttf}
    \newfontfamily\khatt[
        Script=Arabic, Scale=1.3, Path=../../fonts/,
        BoldFont=ae_AlBattar.ttf
    ]{ae_AlBattar.ttf}
}

%------------------------------------------------
% نسخة الطالب / نسخة التصحيح
% لإنتاج نسخة التصحيح: عرّف \withanswers قبل \input لهذا الملف
%------------------------------------------------

\newif\ifcorrection
\ifdefined\withanswers \correctiontrue \else \correctionfalse \fi

\newcommand{\emptybox}{$\square$}
\ifcorrection
    \newcommand{\correctbox}{\textcolor{green!50!black}{$\boxtimes$}}
    \newcommand{\correct}[1]{\textcolor{green!40!black}{\textbf{#1}}}
    \newcommand{\explanation}[1]{%
        \par\vspace{3pt}
        \fcolorbox{gray!60}{gray!8}{\parbox{\dimexpr\linewidth-2\fboxsep-2\fboxrule}{%
        \small\textbf{التعليل:} #1}}}
\else
    \newcommand{\correctbox}{$\square$}
    \newcommand{\correct}[1]{#1}
    \newcommand{\explanation}[1]{}
\fi

\newcommand{\pts}[1]{\textbf{(#1 ن)}}

%------------------------------------------------
% إعداد الصفحة
%------------------------------------------------

\renewcommand{\footrulewidth}{1pt}
\renewcommand{\headrulewidth}{1pt}

\rhead{\khatt المدرسة العليا للأساتذة ببوزريعة}
\lhead{\khatt ملحقة بومرداس}

\cfoot{\thepage}
\lfoot{\khatt الأستاذ: د. بن عبد الرزاق يوسف}

\parindent 0pt
\setlength{\headsep}{15pt}
\arraycolsep=1.4pt

\pagestyle{fancy}

\setlist[enumerate,1]{label=\textbf{السؤال \arabic*:}, leftmargin=*, itemsep=12pt, align=left}

%================================================
% بداية الوثيقة
%================================================

\begin{document}

\begin{center}

{\Large\khatt\textbf{الواجب الثاني}\ifcorrection\ \textbf{-- التصحيح النموذجي}\fi}

\vspace{0.2cm}

{\Large\khatt\textbf{في الذكاء الاصطناعي}}

\end{center}

\rule{\textwidth}{1pt}

\ifcorrection\else
\textbf{الاسم واللقب:} \dotfill \quad \textbf{الفوج:} \makebox[3cm]{\dotfill}

\vspace{0.3cm}
\fi

\begin{center}

{\khatt\textbf{أسئلة متعددة الاختيارات (20 سؤالا -- 15 نقطة)}}

\end{center}

\rule{\textwidth}{1pt}

\textbf{التعليمة:} ضع علامة \textbf{$\times$} في المربع المقابل للإجابة الصحيحة. لكل سؤال إجابة صحيحة واحدة فقط.

\vspace{0.3cm}

\begin{enumerate}

\needspace{9\baselineskip}\item
\textbf{يتعلّم التلميذ الإعراب بقراءة أمثلة مُعرَبة كثيرة مع أستاذه، فيلاحظ القاعدة، ثم يعطيه الأستاذ جملًا جديدة لم يرها من قبل ليتأكد من فهمه. إذا قارنّا هذا المسار بطريقة تعلّم الآلة، فما المرحلة التي تقابل إعطاء التلميذ جملًا جديدة؟} \hfill \pts{0,75}
\begin{itemize}[leftmargin=1.6cm, labelsep=0.3cm, itemsep=2pt, topsep=4pt]
    \item[\emptybox\ \textbf{أ)}] مرحلة التدريب، لأن الآلة تستمر فيها في قراءة أمثلة إضافية، وكلما أعطيناها جملًا جديدة زاد تدريبها، فلا فرق بين جمل التعلّم وجمل التأكد من الفهم.
    \item[\emptybox\ \textbf{ب)}] مرحلة جمع البيانات، لأن الجمل الجديدة تُضاف ببساطة إلى قاعدة الأمثلة، ثم يُعاد بناء النموذج من البداية في كل مرة نريد فيها التحقق من مستواه.
    \item[\emptybox\ \textbf{ج)}] مرحلة النتيجة، لأن الآلة تعرض فيها إجابتها النهائية على المستخدم، ولا حاجة إلى أي تقييم مسبق ما دامت الإجابة قد صدرت عن نموذج مدرَّب بالفعل.
    \item[\correctbox\ \textbf{د)}] \correct{مرحلة الاختبار، لأن النموذج يُجرَّب فيها على أمثلة جديدة لم يرها أثناء التدريب، فنعرف هل اكتشف القاعدة فعلًا أم أنه يحفظ الأمثلة التي تعلّم منها فقط.}
\end{itemize}
\explanation{في تشبيه التلميذ والآلة: قراءة الأمثلة المُعرَبة تقابل بيانات التدريب، وملاحظة القاعدة تقابل اكتشاف الأنماط، وإعراب جمل جديدة يقابل اختبار النموذج ببيانات لم يرها من قبل.}

%------------------------------------------------

\needspace{9\baselineskip}\item
\textbf{طلب أستاذ من طلبته ترتيب المراحل الخمس لتعلّم الآلة كما وردت في الدرس. أي ترتيب صحيح؟} \hfill \pts{0,75}
\begin{itemize}[leftmargin=1.6cm, labelsep=0.3cm, itemsep=2pt, topsep=4pt]
    \item[\emptybox\ \textbf{أ)}] التدريب $\leftarrow$ البيانات $\leftarrow$ الاختبار $\leftarrow$ النموذج $\leftarrow$ النتيجة.
    \item[\emptybox\ \textbf{ب)}] البيانات $\leftarrow$ النموذج $\leftarrow$ الاختبار $\leftarrow$ التدريب $\leftarrow$ النتيجة.
    \item[\correctbox\ \textbf{ج)}] \correct{البيانات $\leftarrow$ التدريب $\leftarrow$ النموذج $\leftarrow$ الاختبار $\leftarrow$ النتيجة.}
    \item[\emptybox\ \textbf{د)}] النموذج $\leftarrow$ البيانات $\leftarrow$ التدريب $\leftarrow$ النتيجة $\leftarrow$ الاختبار.
\end{itemize}
\explanation{تبدأ العملية بجمع البيانات، ثم يتعلّم النظام منها في مرحلة التدريب فينتج النموذج، ثم يُختبر النموذج على أمثلة جديدة، وأخيرًا يقدّم النتيجة.}

%------------------------------------------------

\needspace{9\baselineskip}\item
\textbf{بعد أسابيع من التدريب على آلاف الجمل المصحّحة، قال مهندس: «لقد حصلنا أخيرًا على النموذج». ماذا يقصد بالنموذج حسب ما درسناه؟} \hfill \pts{0,75}
\begin{itemize}[leftmargin=1.6cm, labelsep=0.3cm, itemsep=2pt, topsep=4pt]
    \item[\emptybox\ \textbf{أ)}] مجموعة الجمل المصحّحة نفسها بعد ترتيبها في ملف واحد، لأنها هي المرجع الذي يعود إليه النظام كلما طُرح عليه سؤال.
    \item[\correctbox\ \textbf{ب)}] \correct{البرنامج الذي ينتج عن التدريب، والذي «تعلّم» القاعدة من البيانات فأصبح قادرًا على تطبيقها على أمثلة جديدة.}
    \item[\emptybox\ \textbf{ج)}] الإجابة النهائية التي تظهر على شاشة المستخدم، أي التوقع أو التصنيف الذي يقدّمه النظام بعد استعماله في الواقع.
    \item[\emptybox\ \textbf{د)}] الحاسوب القوي الذي جرى عليه التدريب، لأن جودة النتائج ترتبط أساسًا بقدرة العتاد وسرعة المعالج المستعمل في الحساب.
\end{itemize}
\explanation{النموذج هو ثمرة مرحلة التدريب: برنامج استخرج الأنماط من البيانات. أما البيانات فهي المدخلات، والنتيجة هي ما يقدّمه النموذج لاحقًا.}

%------------------------------------------------

\needspace{9\baselineskip}\item
\textbf{جمع فريق 10000 جملة مُعرَبة، واستعمل 8000 منها ليكتشف النظام قواعد الإعراب، واحتفظ بالباقي جانبًا. ما الاسم الصحيح للجمل الـ 8000؟} \hfill \pts{0,75}
\begin{itemize}[leftmargin=1.6cm, labelsep=0.3cm, itemsep=2pt, topsep=4pt]
    \item[\correctbox\ \textbf{أ)}] \correct{بيانات التدريب، لأنها البيانات التي يتعلّم منها النموذج ويكتشف فيها الأنماط والعلاقات.}
    \item[\emptybox\ \textbf{ب)}] بيانات النتيجة، لأنها الأمثلة التي سيعرضها النظام على المستخدم عندما يطرح عليه سؤالًا.
    \item[\emptybox\ \textbf{ج)}] بيانات احتياطية، لا يستعملها النظام إلا إذا فشل التدريب على الجمل الـ 2000 الأخرى.
    \item[\emptybox\ \textbf{د)}] بيانات الاختبار، لأنها الجزء الأكبر الذي نقيس به جودة النموذج بعد الانتهاء من بنائه.
\end{itemize}
\explanation{بيانات التدريب هي التي يتعلّم منها النموذج، أما الجزء المحتفَظ به جانبًا (2000 جملة) فهو بيانات الاختبار التي تُستعمل بعد التدريب لتقييمه.}

%------------------------------------------------

\needspace{9\baselineskip}\item
\textbf{لدى أستاذة 1000 نص مصنّف حسب النوع الأدبي، وتريد تدريب نموذج على التصنيف ثم تقييمه. أي تقسيم يوافق ما رأيناه في الدرس؟} \hfill \pts{0,75}
\begin{itemize}[leftmargin=1.6cm, labelsep=0.3cm, itemsep=2pt, topsep=4pt]
    \item[\emptybox\ \textbf{أ)}] 200 نص للتدريب و800 نص للاختبار، لأن الاختبار أهم من التدريب في تقييم النموذج.
    \item[\emptybox\ \textbf{ب)}] 500 نص للتدريب، ثم اختبار النموذج على النصوص الـ 500 نفسها للتأكد من حفظها.
    \item[\emptybox\ \textbf{ج)}] استعمال النصوص الـ 1000 كلها للتدريب، ثم إعادة استعمالها نفسها في الاختبار.
    \item[\correctbox\ \textbf{د)}] \correct{800 نص للتدريب و200 نص للاختبار، مع عدم استعمال نصوص الاختبار أثناء التدريب.}
\end{itemize}
\explanation{نقسم البيانات عادة إلى جزء كبير للتدريب (نحو 80\%) وجزء صغير للاختبار (نحو 20\%)، ويجب أن تبقى بيانات الاختبار جديدة على النموذج.}

%------------------------------------------------

\needspace{9\baselineskip}\item
\textbf{اقترح طالب اختبار النموذج بالجمل نفسها التي تدرّب عليها «لأنها مضمونة الإعراب». لماذا يُعدّ هذا الاقتراح غير مناسب؟} \hfill \pts{0,75}
\begin{itemize}[leftmargin=1.6cm, labelsep=0.3cm, itemsep=2pt, topsep=4pt]
    \item[\emptybox\ \textbf{أ)}] لأن الجمل التي تدرّب عليها النموذج تُحذف تلقائيًا من ذاكرته بعد انتهاء التدريب، فلا يمكن إعادة استعمالها تقنيًا في أي مرحلة لاحقة من العمل.
    \item[\emptybox\ \textbf{ب)}] لأن الاختبار بالجمل نفسها يستغرق وقتًا أطول بكثير من الاختبار بجمل جديدة، وهذا هو السبب الوحيد الذي يجعل الخبراء يفضّلون البيانات الجديدة.
    \item[\emptybox\ \textbf{ج)}] لأن جمل التدريب أطول عادة من جمل الاختبار، والمقارنة بين نتائج جمل طويلة وأخرى قصيرة تجعل حساب نسبة النجاح معقدًا وغير ممكن.
    \item[\correctbox\ \textbf{د)}] \correct{لأن النموذج قد يجيب عنها من «الذاكرة» لا من الفهم، كما يفعل التلميذ الذي يُختبر بالأمثلة التي شُرحت في القسم، فلا نعرف هل تعلّم القاعدة فعلًا.}
\end{itemize}
\explanation{نختبر النموذج ببيانات جديدة لنعرف هل فهم حقًا. الاختبار بأمثلة التدريب نفسها يشبه اختبار التلميذ بالجمل التي حفظها، فيعطي صورة مضلِّلة عن مستواه.}

%------------------------------------------------

\needspace{9\baselineskip}\item
\textbf{حصل نموذج لتصحيح الأخطاء الإملائية على نتيجة ممتازة في الجمل التي تدرّب عليها، لكنه أخطأ كثيرًا عندما جرّبته أستاذة على نصوص تلاميذها الجديدة. ما التفسير الأقرب إلى ما درسناه؟} \hfill \pts{0,75}
\begin{itemize}[leftmargin=1.6cm, labelsep=0.3cm, itemsep=2pt, topsep=4pt]
    \item[\emptybox\ \textbf{أ)}] النموذج احتاج إلى وقت إضافي للتكيّف مع الحاسوب الجديد، وستتحسن نتائجه تلقائيًا بعد أيام دون أي تغيير في البيانات أو في التدريب.
    \item[\correctbox\ \textbf{ب)}] \correct{النموذج لم يتعلّم القاعدة تعلّمًا حقيقيًا، بل أتقن الأمثلة التي رآها فقط؛ ولهذا نقيّمه دائمًا ببيانات جديدة لم يرها أثناء التدريب.}
    \item[\emptybox\ \textbf{ج)}] نصوص التلاميذ مكتوبة بخط اليد، والنماذج لا تستطيع أبدًا معالجة أي نص لم يُكتب على الحاسوب، مهما كانت البيانات التي تدرّبت عليها.
    \item[\emptybox\ \textbf{د)}] النتيجة الممتازة في جمل التدريب دليل قاطع على جودة النموذج، فالأخطاء في نصوص التلاميذ سببها حتمًا أخطاء الأستاذة في استعماله.
\end{itemize}
\explanation{النجاح على بيانات التدريب لا يكفي. الاختبار ببيانات جديدة هو الذي يكشف هل اكتشف النموذج الأنماط العامة أم اكتفى بالأمثلة التي رآها.}

%------------------------------------------------

\needspace{9\baselineskip}\item
\textbf{تدرّب نموذج على 500 صورة هواتف و500 صورة حواسيب و100 صورة فقط لألواح رقمية. ماذا نتوقع عندما يرى النموذج لوحًا رقميًا جديدًا؟} \hfill \pts{0,75}
\begin{itemize}[leftmargin=1.6cm, labelsep=0.3cm, itemsep=2pt, topsep=4pt]
    \item[\correctbox\ \textbf{أ)}] \correct{سيخطئ في الغالب أكثر مما يخطئ مع الفئتين الأخريين، لأن صور الألواح قليلة جدًا، فالبيانات غير ممثِّلة لهذه الفئة.}
    \item[\emptybox\ \textbf{ب)}] سيتعرّف عليه بالدقة نفسها التي يتعرّف بها على الهواتف، لأن عدد الصور في كل فئة لا يؤثر في جودة التعلّم.
    \item[\emptybox\ \textbf{ج)}] سيتعرّف عليه بدقة أكبر من غيره، لأن قلّة الصور تجعل النموذج يركّز عليها أكثر ويتذكّر تفاصيلها بشكل أفضل.
    \item[\emptybox\ \textbf{د)}] سيرفض النموذج تصنيف الصورة تمامًا، لأن أي فئة يقل عدد صورها عن 500 تُحذف تلقائيًا من التدريب.
\end{itemize}
\explanation{يتعلّم النموذج ممّا يراه فقط. حين تكون فئة ما قليلة التمثيل في البيانات، يكون تعلّمه لها ناقصًا، فتكثر أخطاؤه فيها.}

%------------------------------------------------

\needspace{9\baselineskip}\item
\textbf{أي وصف يعبّر بدقة عن «البيانات الجيدة» كما عرّفها الدرس؟} \hfill \pts{0,75}
\begin{itemize}[leftmargin=1.6cm, labelsep=0.3cm, itemsep=2pt, topsep=4pt]
    \item[\emptybox\ \textbf{أ)}] بيانات حديثة فقط، لأن الأمثلة القديمة تُربك النموذج ولا فائدة منها أبدًا.
    \item[\emptybox\ \textbf{ب)}] بيانات متشابهة جدًا فيما بينها، حتى يتعلّم النموذج نمطًا واحدًا بسرعة.
    \item[\emptybox\ \textbf{ج)}] بيانات ضخمة جدًا من مصدر واحد، لأن الكمية وحدها تضمن جودة النموذج دائمًا.
    \item[\correctbox\ \textbf{د)}] \correct{بيانات متنوعة وممثِّلة، تغطي مختلف الحالات التي سيواجهها النظام في الواقع.}
\end{itemize}
\explanation{البيانات الجيدة متنوعة وممثِّلة: لا يكفي حجمها، بل يجب أن تغطي الحالات والفئات التي سيتعامل معها النظام عند استعماله.}

%------------------------------------------------

\needspace{9\baselineskip}\item
\textbf{في نقاش حول الانحياز الخوارزمي، طُلب من الطلبة تعريفه بدقة. أي تعريف هو الصحيح؟} \hfill \pts{0,75}
\begin{itemize}[leftmargin=1.6cm, labelsep=0.3cm, itemsep=2pt, topsep=4pt]
    \item[\emptybox\ \textbf{أ)}] خطأ مقصود يضيفه المبرمجون إلى النظام ليوجّه المستخدمين نحو منتجات تجارية معيّنة.
    \item[\correctbox\ \textbf{ب)}] \correct{نتائج غير عادلة أو غير دقيقة يقدّمها النظام، سببها مشكلة في البيانات أو في طريقة تصميمه.}
    \item[\emptybox\ \textbf{ج)}] عطل تقني مؤقت يظهر عند ضعف الاتصال بالإنترنت، ويختفي بمجرد تحسّن جودة الشبكة.
    \item[\emptybox\ \textbf{د)}] ميل النظام إلى تفضيل الإجابات القصيرة على الطويلة، بسبب محدودية سرعة الحاسوب المستعمل.
\end{itemize}
\explanation{الانحياز الخوارزمي هو نتائج غير عادلة أو غير دقيقة، ويعود في الغالب إلى بيانات ناقصة أو غير متوازنة، أو إلى طريقة تصميم النظام.}

%------------------------------------------------

\needspace{9\baselineskip}\item
\textbf{نموذج لتصحيح الأخطاء تدرّب على نصوص صحفية حديثة فقط. عندما أعطته أستاذة بيتًا من الشعر الجاهلي، عدّ بعض الكلمات الفصيحة القديمة أخطاءً. ما سبب ذلك؟} \hfill \pts{0,75}
\begin{itemize}[leftmargin=1.6cm, labelsep=0.3cm, itemsep=2pt, topsep=4pt]
    \item[\emptybox\ \textbf{أ)}] الأستاذة لم تكتب طلبها بصيغة صحيحة، ولو أعادت صياغته لقبل النموذج الكلمات القديمة دون أي حاجة إلى تغيير بيانات التدريب.
    \item[\emptybox\ \textbf{ب)}] الشعر الجاهلي يحتوي فعلًا على أخطاء لغوية كثيرة، والنموذج كان محقًّا في الإشارة إليها لأنه تدرّب على لغة حديثة صحيحة.
    \item[\correctbox\ \textbf{ج)}] \correct{لم يرَ النموذج هذا النوع من النصوص أثناء التدريب، فبياناته لا تمثّل اللغة القديمة، وهذا مثال على الانحياز الناتج عن البيانات.}
    \item[\emptybox\ \textbf{د)}] النموذج يرفض كل نص موزون لأسباب تقنية، فالأوزان الشعرية تُربك الحاسوب مهما كانت البيانات التي تدرّب عليها النموذج.
\end{itemize}
\explanation{هذا مثال الدرس نفسه: النموذج يعكس البيانات التي تعلّم منها. لأنه لم يرَ الشعر القديم، يعدّ كلماته الصحيحة أخطاءً، فالنتيجة منحازة وغير دقيقة.}

%------------------------------------------------

\needspace{9\baselineskip}\item
\textbf{يريد فريق تقليل الانحياز في نموذج لتصنيف النصوص العربية. أي مجموعة من الإجراءات توافق ما درسناه؟} \hfill \pts{0,75}
\begin{itemize}[leftmargin=1.6cm, labelsep=0.3cm, itemsep=2pt, topsep=4pt]
    \item[\correctbox\ \textbf{أ)}] \correct{جمع بيانات متنوعة تمثّل جميع الفئات، والموازنة بين عدد الأمثلة في كل فئة، واختبار النموذج على حالات مختلفة ومراجعة أخطائه.}
    \item[\emptybox\ \textbf{ب)}] تدريب النموذج على مصدر واحد موثوق فقط، لأن تنويع المصادر يُدخل الفوضى، ثم قبول نتائجه مباشرة ما دامت صادرة عن مصدر موثوق.
    \item[\emptybox\ \textbf{ج)}] الإبقاء على البيانات كما هي، والاعتماد على قوة الحاسوب لتصحيح الانحياز تلقائيًا أثناء التدريب، دون الحاجة إلى أي اختبار إضافي.
    \item[\emptybox\ \textbf{د)}] مضاعفة أمثلة الفئة الأكبر فقط لأنها الأهم، وحذف الفئات القليلة من البيانات، ثم الاكتفاء بنسبة النجاح العامة دون مراجعة الأخطاء.
\end{itemize}
\explanation{طرق التقليل من الانحياز في الدرس ثلاث: بيانات متنوعة تمثّل جميع الفئات، وموازنة عدد الأمثلة بين الفئات، واختبار النموذج على حالات مختلفة مع مراجعة أخطائه.}

%------------------------------------------------

\needspace{9\baselineskip}\item
\textbf{رسم أستاذ على السبورة سلسلة توضّح كيف ينشأ الانحياز الخوارزمي. أي سلسلة توافق الدرس؟} \hfill \pts{0,75}
\begin{itemize}[leftmargin=1.6cm, labelsep=0.3cm, itemsep=2pt, topsep=4pt]
    \item[\emptybox\ \textbf{أ)}] نموذج قوي $\leftarrow$ بيانات كثيرة $\leftarrow$ نتائج عادلة دائمًا مهما كان مصدر البيانات أو تنوّعها.
    \item[\emptybox\ \textbf{ب)}] نتائج منحازة $\leftarrow$ بيانات ناقصة $\leftarrow$ نموذج تعلّم صورة ناقصة $\leftarrow$ تدريب جديد على البيانات نفسها.
    \item[\correctbox\ \textbf{ج)}] \correct{بيانات ناقصة أو من فئة واحدة $\leftarrow$ نموذج تعلّم صورة ناقصة $\leftarrow$ نتائج منحازة وغير عادلة.}
    \item[\emptybox\ \textbf{د)}] سرعة الحاسوب $\leftarrow$ حجم الذاكرة $\leftarrow$ نتائج منحازة، دون أي علاقة بالبيانات التي تدرّب عليها.
\end{itemize}
\explanation{ينشأ الانحياز في الغالب من البيانات: بيانات ناقصة أو من فئة واحدة تجعل النموذج يتعلّم صورة ناقصة عن الواقع، فتأتي نتائجه منحازة.}

%------------------------------------------------

\needspace{9\baselineskip}\item
\textbf{قال زميل: «النموذج أعطى هذه النتيجة، إذن هي صحيحة بالضرورة». أي ردّ يوافق ما درسناه؟} \hfill \pts{0,75}
\begin{itemize}[leftmargin=1.6cm, labelsep=0.3cm, itemsep=2pt, topsep=4pt]
    \item[\correctbox\ \textbf{أ)}] \correct{هذا غير كافٍ؛ فالنموذج يعكس البيانات التي تعلّم منها بما فيها من نقص، لذلك تبقى مراجعة نتائجه ضرورية.}
    \item[\emptybox\ \textbf{ب)}] هذا صحيح؛ لأن النموذج لا يمكن أن يخطئ بعد التدريب، فكل أخطائه تُصحَّح أثناء مرحلة جمع البيانات.
    \item[\emptybox\ \textbf{ج)}] هذا صحيح دائمًا ما دام النموذج قد مرّ بمرحلة الاختبار مرة واحدة على الأقل في أي وقت سابق.
    \item[\emptybox\ \textbf{د)}] هذا صحيح إذا كانت النتيجة طويلة ومفصّلة، لأن طول الإجابة دليل على أن النموذج فهم السؤال جيدًا.
\end{itemize}
\explanation{يحذّر الدرس من هذه العبارة بالضبط: لا يكفي أن نقول إن النموذج أعطى النتيجة، فهو يعكس بياناته بما فيها من نقص، والمراجعة مسؤوليتنا.}

%------------------------------------------------

\needspace{9\baselineskip}\item
\textbf{وصف الدرس ما تفعله الآلة أثناء التدريب بأنها «تكتشف الأنماط». ما المقصود بالأنماط هنا؟} \hfill \pts{0,75}
\begin{itemize}[leftmargin=1.6cm, labelsep=0.3cm, itemsep=2pt, topsep=4pt]
    \item[\correctbox\ \textbf{أ)}] \correct{الأشياء التي تتكرر في الأمثلة الكثيرة، مثل القاعدة التي يلاحظها التلميذ في جمل مُعرَبة متعددة.}
    \item[\emptybox\ \textbf{ب)}] الألوان والخطوط المستعملة في عرض الأمثلة، لأن الآلة تحفظ الشكل الخارجي للبيانات قبل محتواها.
    \item[\emptybox\ \textbf{ج)}] القواعد التي يكتبها المبرمج يدويًا داخل البرنامج، قبل أن تبدأ الآلة في قراءة أي مثال من البيانات.
    \item[\emptybox\ \textbf{د)}] الأمثلة النادرة التي لا تتكرر أبدًا، لأن الآلة تهتم بالحالات الاستثنائية أكثر من الحالات الشائعة.
\end{itemize}
\explanation{الأنماط هي الأشياء التي تتكرر في البيانات. كما يلاحظ التلميذ القاعدة من تكرارها في الأمثلة، تكتشف الآلة الأنماط من كثرة الأمثلة.}

%------------------------------------------------

\needspace{9\baselineskip}\item
\textbf{تدرّب نظام للتعرّف على الكلام على تسجيلات لمتحدثين من منطقة واحدة، فصار يخطئ كثيرًا مع لهجات المناطق الأخرى. ما الحل الأنسب؟} \hfill \pts{0,75}
\begin{itemize}[leftmargin=1.6cm, labelsep=0.3cm, itemsep=2pt, topsep=4pt]
    \item[\emptybox\ \textbf{أ)}] الطلب من جميع المستخدمين تقليد لهجة المنطقة الأولى، لأن تغيير البيانات أصعب من تغيير طريقة كلام الناس.
    \item[\correctbox\ \textbf{ب)}] \correct{إضافة تسجيلات متوازنة لمتحدثين من مناطق ولهجات مختلفة، ثم إعادة التدريب واختبار النظام على متحدثين متنوعين.}
    \item[\emptybox\ \textbf{ج)}] إضافة مزيد من التسجيلات من المنطقة الأولى نفسها، لأن كثرة البيانات تكفي وحدها لحل كل المشكلات.
    \item[\emptybox\ \textbf{د)}] الاكتفاء بنسبة النجاح الجيدة على المنطقة الأولى، وعدّ أخطاء المناطق الأخرى حالات استثنائية لا تهمّ.
\end{itemize}
\explanation{المشكلة في عدم تمثيل البيانات للمستخدمين الحقيقيين. الحل هو تنويع البيانات وموازنتها، ثم اختبار النظام على حالات مختلفة.}

%------------------------------------------------

\needspace{9\baselineskip}\item
\textbf{تريد أستاذة تدريب نموذج يميّز بين الشعر والقصة والمقال، فوجدت لديها 900 قصيدة و80 قصة و70 مقالًا. ما أول ما ينبغي فعله قبل التدريب؟} \hfill \pts{0,75}
\begin{itemize}[leftmargin=1.6cm, labelsep=0.3cm, itemsep=2pt, topsep=4pt]
    \item[\emptybox\ \textbf{أ)}] البدء في التدريب فورًا، لأن عدد النصوص الإجمالي كبير بما يكفي لتعلّم الفئات الثلاث.
    \item[\emptybox\ \textbf{ب)}] إضافة مزيد من القصائد، لأن الفئة الكبيرة هي التي تحدّد جودة النموذج كله في النهاية.
    \item[\emptybox\ \textbf{ج)}] حذف القصص والمقالات، وتدريب النموذج على الشعر وحده لأنه الفئة الأكثر حضورًا.
    \item[\correctbox\ \textbf{د)}] \correct{زيادة عدد القصص والمقالات حتى تقترب الفئات الثلاث في العدد، لأن البيانات غير متوازنة.}
\end{itemize}
\explanation{البيانات هنا غير متوازنة: فئة كبيرة جدًا وفئتان صغيرتان. من طرق التقليل من الانحياز الموازنة بين عدد الأمثلة في كل فئة قبل التدريب.}

%------------------------------------------------

\needspace{9\baselineskip}\item
\textbf{اكتشف فريق أن بعض جمل الاختبار أُدرجت خطأً في بيانات التدريب، ثم حصل النموذج على نسبة نجاح مرتفعة جدًا. كيف نفهم هذه النسبة؟} \hfill \pts{0,75}
\begin{itemize}[leftmargin=1.6cm, labelsep=0.3cm, itemsep=2pt, topsep=4pt]
    \item[\emptybox\ \textbf{أ)}] لا تأثير لذلك على النسبة، لأن الاختبار لا يرتبط بالبيانات التي تدرّب عليها النموذج، بل بسرعة الحاسوب فقط.
    \item[\emptybox\ \textbf{ب)}] النسبة أقل من الحقيقة، لأن النموذج يرتبك عندما يرى الجمل مرتين، فيعطي نتائج أسوأ من مستواه الفعلي في الواقع.
    \item[\correctbox\ \textbf{ج)}] \correct{النسبة مبالغ فيها على الأرجح، لأن النموذج رأى جزءًا من جمل الاختبار من قبل، فلم يعد الاختبار يقيس الفهم على بيانات جديدة.}
    \item[\emptybox\ \textbf{د)}] النسبة دقيقة تمامًا، لأن إدراج جمل الاختبار في التدريب يجعل النموذج أذكى، وهذه طريقة مستحبّة لتحسين أي نموذج.
\end{itemize}
\explanation{قيمة الاختبار في أن بياناته جديدة على النموذج. إذا رأى النموذج جمل الاختبار أثناء التدريب، قد يجيب من الذاكرة، فتصير النسبة أعلى من مستواه الحقيقي.}

%------------------------------------------------

\needspace{9\baselineskip}\item
\textbf{جرّبت أستاذة نموذجًا لتصنيف الأخطاء النحوية، ووجدت أنه يخطئ في جمل الفعل المبني للمجهول. ما الخطوة التي تنسجم مع الدرس؟} \hfill \pts{0,75}
\begin{itemize}[leftmargin=1.6cm, labelsep=0.3cm, itemsep=2pt, topsep=4pt]
    \item[\emptybox\ \textbf{أ)}] التوقف عن استعمال النموذج نهائيًا، لأن أي خطأ في حالة واحدة يعني أن كل نتائجه في الحالات الأخرى خاطئة بالضرورة.
    \item[\emptybox\ \textbf{ب)}] تجاهل هذه الأخطاء، لأن النموذج ما دام قد نجح في أغلب الجمل فهو صحيح، ولا فائدة من مراجعة الحالات القليلة.
    \item[\correctbox\ \textbf{ج)}] \correct{مراجعة أخطاء النموذج في هذه الحالة، والتحقق هل كانت أمثلتها قليلة في البيانات، ثم إضافة أمثلة منها واختباره من جديد.}
    \item[\emptybox\ \textbf{د)}] إخفاء جمل المبني للمجهول عن النموذج عند الاستعمال، حتى لا تظهر الأخطاء، دون أي تغيير في البيانات أو في التدريب.
\end{itemize}
\explanation{اختبار النموذج على حالات مختلفة ومراجعة أخطائه من طرق التقليل من الانحياز. الخطأ المتكرر في حالة معيّنة يشير غالبًا إلى نقص في تمثيلها داخل البيانات.}

%------------------------------------------------

\needspace{9\baselineskip}\item
\textbf{طُلب من طالب أن يلخّص في جملة واحدة كيف يتعلّم الذكاء الاصطناعي. أي خلاصة هي الأدق؟} \hfill \pts{0,75}
\begin{itemize}[leftmargin=1.6cm, labelsep=0.3cm, itemsep=2pt, topsep=4pt]
    \item[\emptybox\ \textbf{أ)}] يتبع قواعد يكتبها المبرمج بيده خطوة بخطوة، ولا يعتمد على البيانات إلا لعرض النتائج على المستخدم في المرحلة الأخيرة.
    \item[\correctbox\ \textbf{ب)}] \correct{يتعلّم من بيانات كثيرة يكتشف فيها الأنماط، ثم يُختبر على بيانات جديدة؛ وجودة ما يتعلّمه مرتبطة بتنوّع البيانات وتمثيلها.}
    \item[\emptybox\ \textbf{ج)}] يحفظ كل الأمثلة التي تُعطى له حرفيًا، ثم يبحث فيها عن إجابة مطابقة؛ فكلما زاد عدد الأمثلة زادت دقته دون أي شرط آخر.
    \item[\emptybox\ \textbf{د)}] يتعلّم من سؤال المستخدم وحده في كل مرة، دون بيانات تدريب مسبقة، ولهذا لا يمكن أن تكون نتائجه منحازة أبدًا.
\end{itemize}
\explanation{خلاصة الدرس: تتعلّم الآلة من البيانات كما يتعلّم التلميذ من الأمثلة، وتُختبر ببيانات جديدة، والبيانات الناقصة تؤدي إلى الانحياز.}

\end{enumerate}

\vfill
\begin{center}
{\khatt بالتوفيق}
\end{center}

\end{document}
